\chapter{Background and Related Work}
\label{ch:background}

\section{Alpha research in plain terms}
\label{sec:bg-primer}

This report uses a small vocabulary that is standard in quantitative
finance and opaque outside it; this section defines it once, so that
the rest of the report does not presuppose a finance background.

An \emph{alpha factor} is a rule that assigns every asset a score at
every point in time, constructed so that higher-scored assets should
subsequently outperform lower-scored ones. The prediction is about
\emph{relative} performance within a fixed universe, not about price
levels: the object of interest is the \emph{cross-section} --- all
assets observed at one timestamp --- and whether today's ranking of
scores anticipates the ranking of returns realised over the following
days. A factor that scores an asset from that asset's own history
alone (five-day price reversal, abnormal trading volume) is an
\emph{island factor} in this report's terminology; a \emph{relational
factor} may also read the asset's neighbours in a graph.

Two measurements recur on every page of \cref{ch:equity,ch:energy}.
The \emph{information coefficient} (IC) is the correlation between the
scores issued at a timestamp and the returns realised afterwards; this
report uses the rank form (\emph{rank-IC}, a Spearman correlation),
computed cross-sectionally at each timestamp and averaged. A rank-IC
of 1 would be a perfect ordering; 0 is no predictive skill at all.
Real factors are weak --- the best out-of-sample composite in this
report reaches a rank-IC of about 0.015 (\cref{ch:equity}) --- and
factors of this strength are useful only in combination, which is why
the field builds them in large numbers. The \emph{Sharpe ratio}
evaluates a factor as a trading strategy instead: buy the top-ranked
assets, sell short the bottom-ranked ones (a \emph{long-short}
portfolio, whose exposure to the overall market direction largely
cancels), and divide the strategy's mean return by its volatility,
annualised. The two numbers can disagree, because ranking well says
nothing about the magnitudes of wins and losses; \cref{ch:energy}
turns exactly this disagreement into one of the report's main
findings.

``Alpha'' itself, in the industry idiom this report inherits from
Kakushadze's catalogue of practitioner formulas
\cite{kakushadze2016alphas}, names nothing deeper than a signal
expected to carry predictive information beyond passive market
exposure. No claim below depends on a stricter definition.

\section{Cross-sectional alpha research}
\label{sec:bg-alpha}

The equity track's island factors follow the formulaic-alpha idiom that
Kakushadze documented from industry practice: short arithmetic
expressions over price and volume, evaluated per stock and ranked across
the cross-section each day \cite{kakushadze2016alphas}. Individually
weak and mutually correlated, such alphas are meant to be combined, and
the interesting research questions live in the combination step. Machine
learning has a measured track record here: Gu, Kelly and Xiu showed on
six decades of US data that flexible models --- trees and shallow
networks in particular --- roughly double the predictive $R^2$ of linear
benchmarks for cross-sectional returns, driven largely by interaction
effects that linear factor models cannot express
\cite{gu2020empirical}. The same literature is emphatic about how easily
such gains evaporate under sloppy evaluation. L\'opez de Prado catalogues
the failure modes --- label overlap, missing embargoes, selection on
test data, backtest overfitting --- and the purged, embargoed validation
splits he advocates \cite{lopezdeprado2018afml} are the direct ancestors
of the split protocol in \cref{sec:splits}. This report treats that
methodological literature as load-bearing: two of the energy track's
three artifacts (\cref{sec:en-artifact}) are textbook instances of
exactly the pathologies it names.

\section{Graph neural networks and attention}
\label{sec:bg-gnn}

Graph neural networks compute node representations by aggregating
information from graph neighbourhoods. In the graph convolutional
network of Kipf and Welling, the aggregation weights are fixed by the
graph's degree structure \cite{kipf2017gcn}; graph attention networks
replace them with weights computed from the node features themselves, so
the model learns which neighbours matter for each node
\cite{velickovic2018gat}. Brody, Alon and Yahav later showed that the
original formulation computes a restricted, effectively static form of
attention and proposed GATv2, which reorders the operations to make the
attention genuinely conditional on the query node
\cite{brody2022gatv2}. Both operators are available off the shelf in
PyTorch Geometric \cite{fey2019pyg}; the alpha track uses the original
operator, the forecasting work the GATv2 variant
(\cref{sec:training}).

Applying relational models to stock prediction is an established idea.
Feng et al.\ combine a temporal encoder with relational graphs over
company links and learn to rank stocks by return
\cite{feng2019stockranking}, and much subsequent work follows the same
pattern: propose an architecture, report that it beats non-relational
baselines. This project asks a deliberately narrower question. Rather
than demonstrating that some graph-augmented architecture outperforms,
it holds the information set fixed --- the same island alphas every
baseline consumes --- and measures the marginal value of each
ingredient separately: the graph via an unlearned uniform-propagation
anchor, and learned attention via the gap between the GAT and that
anchor (\cref{sec:anchors}). The price of the narrower question is
weaker headline numbers; the payoff is that when the equity track
reports a positive attention margin in thirty of thirty runs
(\cref{ch:equity}), the claim isolates a mechanism rather than a
model.

\section{Electricity price forecasting}
\label{sec:bg-epf}

Day-ahead electricity prices are set in a daily auction and driven by
fundamentals --- load, wind and solar in-feed, fuel prices --- with
properties equity researchers rarely face: strong diurnal seasonality,
negative prices, and fat-tailed scarcity spikes. Weron's review maps the
field's method families and stresses the role of these fundamentals and
of rigorous out-of-sample evaluation \cite{weron2014epf}; Lago et al.\
sharpen the evaluation point for the deep-learning era, offering an
open benchmark and best practices --- among them, always comparing
against naive and well-tuned simple baselines --- precisely because the
literature is rich in models that beat only weak straw men
\cite{lago2021epf}. The forecast-skill ladder of \cref{sec:skill}, with
persistence pinned at zero skill and a ridge on fundamentals as the
serious no-graph baseline, is this project's implementation of that
advice. Its energy-specific contribution is not a new forecasting
architecture but a controlled measurement of one factor the standard
benchmarks do not isolate: whether the physical interconnector topology
itself adds forecast skill, at the node level and --- where the answer
turns out to be strongest --- for cross-border spreads
(\cref{sec:en-edge}).

\section{Where this report sits}
\label{sec:bg-position}

Against this background, three commitments distinguish this report from
the related work above, which rarely combines them. Every relational
claim is measured against matched no-learning anchors on identical
inputs, so a positive result isolates an ingredient rather than
crediting a whole architecture. One GAT kernel and one evaluation
harness are carried across two structurally different graphs --- an
estimated correlation backbone over equities and a physical transmission
network over bidding zones --- with the negative controls travelling to
both. And the honesty standard of the methodological literature is
applied without exception: negative results (a failed value-added gate,
a money-losing energy strategy, a null congestion hypothesis) are
reported as findings, because the diagnostic arc that produced them is
the part of this work most likely to help a reader.
